\chapter{研究设计与测量}
\label{ch:methodology}
\chaptermark{研究设计与测量方法}
评估需要把系统干预与可观测的成本或质量变化联系起来。本文结合配对执行消融、
留出查询上的预算选择，以及各系统独立调优后的前沿比较。这些实验回答不同问题：
更快的内核、有效的校准策略和有竞争力的搜索系统，各自需要不同的对照。
本章定义这些对照，以及解释结果所需的证据；第4--6章介绍机制，第7章报告实测结果。

\section{将研究问题转化为受控比较}
\label{sec:3-1}
\Cref{tab:3-design}明确比较单位，以及哪些观察会挑战各项研究主张。
一个\emph{实验单元（cell）}固定数据集、索引配置和工作负载设置；
各实验臂是在该单元内替换的实现或预算。配对执行测量固定索引状态、查询、
probe数、精排预算和打分表示，除非这些正是指定的干预变量。
配对比值为$Q_{\rm treatment}/Q_{\rm control}$；几何平均对指定单元的
乘性比值作汇总，各单元等权。

\par\medskip\noindent\begin{minipage}{\linewidth}
\centering\captionsetup{hypcap=false}
\captionof{table}{把研究问题转化为可检验的比较。容量观察与配对执行干预需要不同对照。}
\label{tab:3-design}
\footnotesize\setlength{\tabcolsep}{3pt}
\begin{tabularx}{\linewidth}{@{}p{.11\linewidth}XXX@{}}
\toprule
问题 & 干预／比较 & 固定或单独核算的条件 & 指标与反证 \\
\midrule
RQ1：构建与容量 & GPU构建；缓存与冷构建生命周期；压缩与原始存储
& $N,d$、表示、路由、构建计时边界
& train/add时间与常驻字节；构建失败或模型与分配差额限定可行性 \\
RQ2：执行 & LUT分组、word访问、精排lane、分组扫描、共享界与probe选择
& 配对状态、工作负载、分数及参考输出语义
& 墙钟QPS、阶段时间和输出检查；变慢或非并列差异挑战机制 \\
RQ3：预算 & 抽样／全池规则，对照固定64、100、200、枚举与GT oracle
& 同一索引与路由；校准和留出查询位置互斥
& 留出recall损失、选中臂QPS与校准成本；越阈挑战可靠性 \\
RQ4：工作区间 & 独立调优的JHQ、CAGRA、IVF-PQ和IVF-RaBitQ
& 导出向量、GT、$k$、设备和主机到主机边界；显存另报
& 同recall QPS与覆盖；其他recall处更强的竞争者挑战宽泛优势 \\
\bottomrule
\end{tabularx}
\end{minipage}\par\medskip

研究设计区分机制的实现效果与其原因解释。配对word访问消融可以确定该访问
路径在所测流水线内的吞吐变化；把变化归因于cache miss或occupancy，
还需要硬件计数器。跨系统前沿则比较表示、算法与实现的整体选择，无法单独
识别某个内核的贡献。\Cref{sec:7-6}提供集成调度比较，避免把不同基线上的
组件提速比直接相乘。

这些研究问题组织的是迭代开发记录，并非预注册方案：E1发现recall覆盖不足后
扩展了网格，工作负载精排需求的解释也来自预算扫描。低recall失利、高recall
反转和失败方案同样进入评估。与量化参考的输出一致性、相对原始向量GT的recall，
按定义就是不同的量；预算规则使用前者，留出评估检验后者。

\section{数据集、查询、真值与划分}
\label{sec:3-2}
\Cref{tab:3-data}列出新版预算、LUT与规模实验使用的六份数据导出；
海报E1比较使用前四行。Vogue是时尚图片嵌入，arXiv是摘要嵌入；
两个OpenAI数据集来自Qdrant的DBpedia导出
\cite{data_vogue,data_arxiv,data_openai1536,data_openai3072}。
BGE使用意大利语Wikipedia配置，Stella使用TREC24 BioGen语料导出
\cite{data_bge,data_stella}。规模来自准备后的实测文件，而不是当前上游数据卡。

\par\medskip\noindent\begin{minipage}{\linewidth}
\centering\captionsetup{hypcap=false}
\captionof{table}[数据集与训练配置]{数据与JHQ训练清单。$n_q$为实际查询数；$n_t=39n_{\rm list}$是请求的
粗聚类训练样本数。E1使用前四行；全部使用平方L2及$k=10$。}
\label{tab:3-data}
\small\setlength{\tabcolsep}{4pt}
% Generated by scripts/audit_ch3.py from the budget shape logs.
\begin{tabular*}{\linewidth}{@{\extracolsep{\fill}}lrrrrr@{}}
\toprule
数据集 & $N$ & $d$ & $n_q$ & $n_{\rm list}$ & $n_t$ \\
\midrule
Vogue & 932,328 & 768 & 1,000 & 4,096 & 159,744 \\
arXiv & 2,253,000 & 768 & 1,000 & 8,192 & 319,488 \\
OpenAI-1536 & 999,000 & 1,536 & 1,000 & 8,192 & 319,488 \\
OpenAI-3072 & 999,000 & 3,072 & 1,000 & 8,192 & 319,488 \\
BGE-M3 & 10,091,524 & 1,024 & 1,000 & 32,768 & 1,277,952 \\
Stella & 17,776,615 & 1,024 & 999 & 32,768 & 1,277,952 \\
\bottomrule
\end{tabular*}

\end{minipage}\par\medskip

准备源码规定L2归一化，以及带种子的查询抽取与数据库移除。
流式OpenAI/BGE/Stella转换器先删除非有限行，对范数大于$10^{-12}$的向量
归一化，并用seed 42无放回选择查询位置。Stella查询来自corpus抽样，
没有使用单独的65条官方benchmark查询；较短的实际导出得到999条查询。
这些实验评估向量近邻，不评估带相关性标签的检索。
两个OpenAI导出采用相同位置抽样规则；补证向量核对支持全部查询及10,299行base
样本的按行对应，但未归档原始source-ID连接（\cref{sec:7-precision,sec:7-2}）。

同一单元的所有方法读取同一查询和GT文件。GT通过对导出FP32向量进行
穷举平方L2搜索得到；流式转换器在数据库分块之间维护top-$k$。
评估沿用存储的GT顺序及其浮点并列选择。
设$\mathcal T_k(q)$为GT前$k$个ID，$\mathcal R_k(q)$为返回ID集合：
\begin{equation}
\Recall@k=\frac{1}{|\mathcal Q|k}
\sum_{q\in\mathcal Q}|\mathcal R_k(q)\cap\mathcal T_k(q)|.
\label{eq:3-recall}
\end{equation}
Vogue/arXiv的GT行宽为100，其余四份为20，评分只比较前10项。
返回重复ID不会重复得分；越界、重复和跨调用变化分别记录。

JHQ采用$M=d/8$、$B=B_r=8$、8次粗聚类迭代，并请求100,000条残差训练样本。
搜索扫描复用\texttt{arms6\_cache}训练状态，重新编码数据库；
冷构建从没有该状态开始。源码从数据库抽取训练样本，与校准查询分开。
预算实验以偶数查询位置构成500条校准池，奇数位置构成500条留出池，
Stella留出499条。每个单元在$S=32,64,128$下各抽64次，FULL使用全部500条。
GT只参与离线oracle和留出质量评估，可部署的一致性规则使用搜索输出
（\cref{sec:6-2,sec:a-budget-contract}）。

准备源码和形状日志支持上述协议，但历史向量文件与训练状态的完整哈希缺失。
证据登记区分已检查的准备流程与每次运行数据的不可变身份。

\section{测量边界、波动与显存核算}
\label{sec:3-3}
主比较从主机内存中的FP32查询开始，到流完成后top-10 ID与分数回到主机为止
（\cref{fig:3-boundary}）。区间包含pinned staging、H2D、查询预处理、搜索和D2H。
CAGRA int8查询量化计入其中，量化器训练和索引构建在区间外。
资源句柄与输出缓冲复用。JHQ回传到pageable数组，基线回传到pinned数组；
CAGRA/PQ经Python API，JHQ/RaBitQ经C++。解释微小差距时保留这些接口差异。

\par\medskip\noindent\begin{minipage}{\linewidth}
\centering\captionsetup{hypcap=false}
\resizebox{\linewidth}{!}{\begin{tikzpicture}[
 font=\small\linespread{1}\selectfont,
 box/.style={draw=ThesisNavy!65,rounded corners=2pt,align=center,
 text width=2.55cm,minimum height=1.55cm,inner sep=4pt},
 arr/.style={-{Latex[length=2mm]},draw=ThesisNavy!80,thick}]
\node[box] (host) at (0,0) {主机FP32查询\\Pinned staging};
\node[box,fill=ThesisNavy!7] (upload) at (3.5,0) {H2D传输\\查询预处理};
\node[box,fill=ThesisNavy!7] (gpu) at (7,0) {GPU搜索\\选择／精排};
\node[box] (result) at (10.5,0) {D2H ID与分数\\等待流完成};
\draw[arr] (host) -- (upload);
\draw[arr] (upload) -- (gpu);
\draw[arr] (gpu) -- (result);
\draw[ThesisNavy,thick] (-1.4,1.15) -- (-1.4,1.45) -- (11.9,1.45) -- (11.9,1.15);
\node[fill=white,inner sep=3pt] at (5.25,1.45) {主口径：主机到主机计时区间};
\node[box,text width=5.7cm,minimum height=1.3cm] at (1.65,-2.35)
 {搜索计时之外\\加载数据／训练或加载模型／add／预热};
\node[box,text width=5.7cm,minimum height=1.3cm] at (8.65,-2.35)
 {独立策略成本\\预算校准；基于模型的复用回本};
\draw[densely dashed,ThesisNavy!60] (-1.4,-1.1) -- (11.9,-1.1);
\end{tikzpicture}
}
\captionof{figure}[主机与设备计时边界]{计时边界。主机区间包含传输和查询相关计算；设备阶段profile解释其组成。
单独profile与同步可能改变执行，因此阶段时间之和不能替代端到端延迟。}
\label{fig:3-boundary}
\end{minipage}\par\medskip

第$r$轮有$m_r$次计时调用，每次回答$b$条查询：
\begin{equation}
Q_r=\frac{m_rb}{\sum_{i=1}^{m_r}t_{ri}},\qquad
Q=\operatorname{median}_r Q_r,
\label{eq:3-throughput}
\end{equation}
其中$t$以秒计。E1采用$b=1{,}000$和5轮测量。
JHQ每个变体每轮先做1次不计时调用，再通常计时24次；基线先做3次初始预热，
每轮通常计时24次，对判定较慢的点减少次数，并保留实际次数。
预热使JIT和graph capture不进入稳态计时。batch调用的中位／p95延迟描述重复
batch调用，不是逐查询或并发请求的延迟分布。

各轮吞吐的中位数与单次调用延迟中位数的倒数不同。在可得时同时保留二者、
各轮最小／最大值及基线逐调用轨迹。有效慢调用保留在主估计中。
各轮共享训练索引，预算抽样共享查询池：其次数衡量计时波动或抽样敏感性，
不等于独立训练重复数。微小配对差异结合离散度和一致性解释，
不对扫描单元附加依赖独立性假设的显著性主张。

主环境为RTX 5090，可报告显存32,607\,MiB，驱动595.71.05，
CUDA toolkit 13.0和cuVS 26.08.01。运行获取协作GPU锁并保存环境快照；
25核容器配额内使用16个OpenMP线程。E1记录功率上限575\,W，
快照并不证明锁频或恒定宿主负载。JHQ Release的compute\_80 PTX由驱动JIT\cite{cuda_platform}，
库基线使用安装的构建。辅助设备常驻计时与阶段profile分别标注，
因为JHQ主API接受主机查询。

显存核算区分持久编码数组、模型／IVF元数据、查询工作区和运行时／context。
进程used样本包含索引以外的分配；构建前后差值不包含之后创建的内存池。
E1中减context后的JHQ总量与CAGRA索引差值保留各自定义。
组件模型解释常驻总量并显示未归因余量（\cref{sec:4-4,sec:7-7}），
不能代替等显存性能比较；补证以20\,ms设备采样另测构建峰值。冷train/add计时同样不含输入加载、
此前的对象构造与主机模型设置。

\paragraph{补证的测量边界。}
补证实验新增的检查见\cref{sec:a-supplement}。硬件型号及报告驱动／toolkit与归档一致，
但计时属于另一会话。算术模式及缓存模型下的精度干预见\cref{sec:4-2,sec:7-precision}；
历史基线比较保留会话区别。

\section{基线、调参网格与前沿构造}
\label{sec:3-4}
竞争方法涵盖图遍历（CAGRA fp32/int8）、乘积量化IVF与GPU IVF-RaBitQ
\cite{cagra,pq,ivfrabitq}，FP32 GPU暴力搜索提供穷举参考。
各方法扫描自己的实现参数（\cref{tab:3-grids}），索引配置的存储量和重建精度
可以不同。因此主前沿衡量这些网格内指定质量下的可达速度，另行报告显存。

\par\medskip\noindent\begin{minipage}{\linewidth}
\centering\captionsetup{hypcap=false}
\captionof{table}{E1主要调参维度。逐点参数和覆盖扩展保留在原始记录中；
各方法使用不同的有限网格。}
\label{tab:3-grids}
\footnotesize\setlength{\tabcolsep}{4pt}
\begin{tabularx}{\linewidth}{@{}p{.22\linewidth}X@{}}
\toprule
方法 & 记录的扫描维度 \\
\midrule
JHQ & probe与$\alpha$；初始probe为64/128/256，预算为工作负载基准的
$0.5/1/2$倍；覆盖扩展最高至probe 1,024，并增加更大／更小预算 \\
CAGRA fp32/int8 & 图度32/64；search width为1/2/4/8；
中间top-$k$容量32--512 \\
IVF-PQ & 8-bit子码；$d=768$时码长192/384/768字节，
更高维度时384/768/1,536；fp16/fp32查表设置；probe 16--512 \\
IVF-RaBitQ & 每维4/8/9 bit、支持的搜索模式、probe 16--1,024；
OpenAI另有4,096-list索引 \\
GPU暴力搜索 & 穷举FP32平方L2，采用相同主机计时边界 \\
\bottomrule
\end{tabularx}
\end{minipage}\par\medskip

前沿保留recall与QPS上不被支配的点：若$R_a\geq R_b$且$Q_a\geq Q_b$，
至少一项严格更好，则$a$支配$b$。包络可以由不同索引配置组成。
对目标recall $r$，实测估计为
\begin{equation}
Q_{\rm meas}(r)=\max\{Q_i:R_i\geq r\}.
\label{eq:3-measured-frontier}
\end{equation}
没有合格样点时无定义。插值估计使用相邻前沿点$(R_a,Q_a)$与$(R_b,Q_b)$：
\begin{equation}
\log Q_{\rm int}(r)=(1-\lambda)\log Q_a+\lambda\log Q_b,
\qquad \lambda=\frac{r-R_a}{R_b-R_a}.
\label{eq:3-interpolated-frontier}
\end{equation}
归档parser先把与目标相差不超过$5\times10^{-4}$的前沿点视为直接匹配，
包括接近端点的情形；否则只在实测范围内部插值，不对线段外推。
差距接近时报告两种估计：前者选实际达到阈值的配置，后者降低网格间距不等
的影响；二者都不构成等存储比较。

parser把失败与非法输出诊断保留在有效点前沿之外。仅有跨调用ID变化不会导致
删点，原始记录区分集合变化与顺序变化。隔离后的大型基线重跑替代受污染构建，
Stella因batch 1,000被跳过的运行不贡献搜索点，决策见\cref{sec:a-e1-protocol}。
论文从全部明确列出的基线族重新构造前沿，包含CAGRA int8及不利工作区间
（\cref{fig:7-frontiers}）。历史ADC与海报图保留各自版本及测量边界。

\section{带版本的主张与证据登记}
\label{sec:3-5}
各实验族记录实现版本、索引与查询配置、计时边界及测量输出。配对比较复用相同
训练状态；版本、状态或精度的变化分别构成独立实验族。历史可执行文件和训练状态
指纹仍不完整，限制了部分运行的精确复现。\Cref{app:reproducibility}汇总可获得的
环境信息与实验协议。

\section{证据强度与有效性控制}
\label{sec:3-6}
证明在明确前提下建立性质，独立参考检验有限样例中的实现一致性。
配对基准测量固定工作负载内的干预效果，系统测量定位被测设备与网格中的竞争区间。
字节与回本模型保留各自核算假设，解释候选需要另行检验。
精确primary选择针对实际算出的primary分数，原始向量recall还受路由和重建影响。


共享索引依赖、有限池、网格选择、接口不对称及单设备范围构成上述对照设计的动机。
总体与排序假设见\cref{sec:6-2,sec:6-3}，\cref{sec:8-3}统一讨论其对解释和部署的限制。

